\documentclass[anon]{nesy2026}

\usepackage{booktabs}
\usepackage{microtype}

% jmlr.cls already loads amsmath, amssymb, natbib; uses ntheorem
\theorembodyfont{\upshape}
\theoremheaderfont{\bfseries}
\theorempostheader{.}
\theoremsep{\newline}
\newtheorem{prop}{Proposition}

\title[TLTS-Compilation]{TLTS-Compilation: A Neurosymbolic Framework
  for Type-Safe and Verifiable Transformers}

\begin{document}

\maketitle

\begin{abstract}
Verifying that a language model stayed within the rules of a
structured domain---that it did not invent an absent relation or
skip a required step---is currently faith-based. We frame the problem
as compiling a typed labeled transition system (TLTS) into a
transformer, and show that two recent threads---typed attention and
program-compiled transformers---are the same construction, differing
only in whether the transition relation is deterministic. The
framework names three points in the inference pipeline where the
rule can be enforced (in-FFN gates, pre-decoder logit masks, post-hoc
audit), each with a distinct cost profile. A synthetic harness
confirms its predictions, including a non-obvious one:
attention-layer reachability masking is not enough for trajectory
soundness; the decoder must also check direct-edge admissibility. We
deliver a JSON audit-certificate format that a third party can
re-check without model weights, and argue it should become a
publication norm for compiled-transformer work. Trained-model
evaluation is future work.
\end{abstract}

\section{Introduction}
\label{sec:intro}

A deployed language model is asked to generate inside a structured
domain---an ontology, a protocol, a program semantics. Today, nothing
in the forward pass forces it to stay there: any token may attend to
any other, and any next-token may be sampled, regardless of whether
the implied move is admissible in the domain. Two recent research
threads address this from opposite ends. From the symbolic side,
work on type-safe attention~\citep{constrained_decoding} gates
attention by a domain ontology---an \emph{olog}, in the sense of
\citet{spivakkent2012ologs}---and constrains generation to outputs
that admit a proof of admissibility. From the neural side, work on compiled
transformers \citep{lindner2023tracr,perez2021turing} compiles
deterministic programs directly into transformer weights, so one
forward pass executes one program step.

These threads look unrelated---one is about typed knowledge graphs,
the other about machine code. Both compile a discrete, typed
transition system into transformer mechanics so the forward pass
implements the system's semantics, but each line states the pattern
in its own vocabulary. We propose stating it once.

We take that object to be the \emph{typed labeled transition
system} (TLTS). Compiling a TLTS into a transformer is a functor
from its path category into a category of compiled-transformer
specifications: each legal path of state transitions maps to a
forward-pass behavior, and gluing paths end-to-end corresponds to
running the computation in sequence (composition is preserved). The
paper develops the framework, identifies three
distinct architectural loci where the transition relation can be
enforced, and validates the resulting predictions empirically. Our
contributions:

\begin{enumerate}
\item A categorical formalization of TLTS-compilation as a functor
  from the path category of a TLTS to a category of compiled
  transformer behaviors (Section~\ref{sec:formal}).
\item A subsumption: procedural (program-compilation) and
  ontological (Olog-compilation) constructions are the deterministic
  (functional-$\delta$) and nondeterministic (relational-$\delta$)
  specializations of the same recipe (Section~\ref{sec:subsumption}).
\item Three enforcement loci (in-FFN, pre-decoder, post-hoc) with
  different cost regimes (Section~\ref{sec:loci}).
\item A prediction we confirm empirically: attention-layer
  reachability masking is insufficient for trajectory soundness; the
  decoder must also check direct-edge admissibility
  (Section~\ref{sec:expts}).
\item A post-hoc verification protocol with an audit-certificate
  reference implementation that admits independent re-checking
  without model internals (Section~\ref{sec:verif}).
\end{enumerate}

\section{TLTS and the Path Category}
\label{sec:formal}

\paragraph{What we need from the math.} The transformer is supposed
to follow a domain rulebook: a finite, declared list of state shapes,
allowed moves between them, and the labels those moves carry. The
smallest object that captures all three is a labeled transition
system; adding \emph{types} to the states lets the rulebook
constrain by state shape rather than by state identity. The rest of
this section attaches one such object to a transformer's forward
pass, so that one forward step on a token corresponds to one move in
the rulebook. The properties we care about downstream---soundness
(Section~\ref{sec:formal}), enforcement loci (Section~\ref{sec:loci}),
and verifiability (Section~\ref{sec:verif})---are read off this
correspondence.

A \textbf{typed labeled transition system} (TLTS) is a tuple $M =
(T, L, \delta, t_0)$ where $T$ is a set of types (state-shape
classes), $L$ is a set of labels (typed tokens), $\delta \subseteq T
\times L \times T$ is a typed transition relation, and $t_0 \in T$
is the initial type. We write $t_1 \stackrel{\ell}{\rightarrow} t_2$
for $(t_1, \ell, t_2) \in \delta$. $M$ is \emph{functional} if
$\delta$ is a partial function $T \times L \rightharpoonup T$; it is
\emph{relational} otherwise. We have
$\mathsf{TLTS}_{\mathrm{func}} \subsetneq
\mathsf{TLTS}_{\mathrm{rel}}$.

\paragraph{Path category.} Every TLTS $M$ induces a category
$\mathcal{C}_M$ whose objects are types in $T$ and whose morphisms
$t_1 \to t_2$ are finite label sequences $(\ell_1, \ldots, \ell_n)$
realizing some trajectory in $\delta$ from $t_1$ to $t_2$.
Composition is concatenation; identities are empty sequences. This
is the free category on the directed multigraph of $\delta$. When
$M$ is functional, $\mathcal{C}_M$ is \emph{thin} (at most one
morphism per ordered pair); when $M$ is relational, parallel
morphisms exist. For a domain Olog $\mathcal{O}$
\citep{spivak2014category}, $\mathcal{C}_\mathcal{O}$ is the
classifying category in the standard sense.

\paragraph{Compiled-transformer specifications.} Fix a transformer
architecture (layers, heads, hidden dimension $d$). A
\emph{compiled-transformer specification} is a triple
$(E, A, \Phi)$ where $E: T \to \mathbb{R}^d$ is a type embedding
(images separated under the residual-stream metric), $A \in
\{0,1\}^{T \times T}$ is an attention admissibility mask, and $\Phi$
is a per-label gate-and-transition assignment $(g_\ell, \tau_\ell)$
realized in FFN rows. Specifications form a
category $\mathsf{CompT}$ under embedding-preserving simulation.

\paragraph{The compilation functor.} A \emph{TLTS-compilation} of
$M$ is a functor $F_M: \mathcal{C}_M \to \mathsf{CompT}$ assembled
from a recipe: choose $E$ injective on $T$; set $A(t_q, t_k) = 1$
iff $t_k$ is reachable from $t_q$ in $\delta$; realize $\delta$ at
one or more of the three enforcement loci of
Section~\ref{sec:loci}. We say a compiled transformer $C$
\emph{realizes} $F_M$ if its forward pass on label sequence
$\ell_1, \ldots, \ell_n$ produces residuals $E(t_0), E(t_1), \ldots,
E(t_n)$ with $(t_{i-1}, \ell_i, t_i) \in \delta$ at each step.

\begin{prop}[Trajectory-soundness]\label{prop:sound}
Let $C$ realize $F_M$ with $E$ injective and per-label gates $g_\ell$
sharp. Then for every input token sequence on which $C$ produces a
residual trace, that trace decodes to a valid trajectory of $M$.
\end{prop}

The proof is direct: the attention mask $A$ blocks information flow
that violates type reachability; gate sharpness ensures exclusive
$\ell$-row activation; injectivity of $E$ guarantees unique type
recovery. Composition along the forward pass thus traces an
admissible path in $\delta$.

\section{Three Enforcement Loci}
\label{sec:loci}

$\Phi$ is left open because the transition relation can be enforced
at three points in the inference pipeline. The three are not
exclusive; production systems combine them. Costs differ:

\begin{table}[h]
\centering
\small
\begin{tabular}{@{}llll@{}}
\toprule
Locus & Where $\delta$ checked & Per-step cost & Architectural cost \\
\midrule
\textbf{In-FFN} & forward pass; FFN row gates & $O(1)$ per row & FFN re-compile \\
\textbf{Pre-decoder} & sampling; logit mask & $O(|L|)$ mask & none (drop-in) \\
\textbf{Post-hoc} & after gen.; trace audit & $O(n)$ for length-$n$ & none on model \\
\bottomrule
\end{tabular}
\caption{Three loci for $\delta$-enforcement.}
\label{tab:loci}
\end{table}

The \textbf{in-FFN} locus is what compiled-transformer constructions
\citep{lindner2023tracr} use; gates \emph{are} the enforcement. It
is the strictest: the model cannot violate $\delta$ regardless of
input, at the cost of constructing or fine-tuning weights for the
specific TLTS.

The \textbf{pre-decoder} locus is the constrained-decoding setting
\citep{willard2023outlines, geng2023grammar}. At each step, the
admissible label set for the current type-state is computed from
$\delta$ and logits over $L$ are masked accordingly. The model is
unchanged; enforcement happens at the sampling boundary.

The \textbf{post-hoc} locus does not enforce $\delta$ but audits
that it was respected. Used alone, it implies a verify-and-retry
loop that wastes compute; used in combination with one of the first
two loci, the verifier becomes a monitoring artifact rather than a
flow-control gate (Section~\ref{sec:verif}).

\section{Subsumption: Procedural is Functional-$\delta$ Compilation}
\label{sec:subsumption}

Procedural and ontological transformer constructions instantiate the
same compilation recipe with different transition systems. A
deterministic virtual machine has $T = $ machine-state shapes,
$L = $ opcodes, $\delta = $ deterministic operational semantics.
Compiled-transformer constructions \citep{lindner2023tracr} bake
$\delta$ into FFN gates (the in-FFN locus); a deterministic
attention-key trick---in 2D embedding $(\text{step}, \text{value})$
with query direction $[1, 0]$, the most-recent-write is recovered as
the argmax of a linear objective, which lies on the convex hull of
the key set, giving $O(\log h)$ retrieval with hull-based
indexing---substitutes for soft attention.

An Olog $\mathcal{O}$ is a TLTS with $T = $ Olog types, $L = $
morphism labels, $\delta = $ Olog edges plus equational
identifications. The ontological attention layer used here
\citep{constrained_decoding} realizes $A$ via reachability masking
over $\mathcal{C}_\mathcal{O}$, with the categorical substrate from
\citet{spivakkent2012ologs}. Proof-guided generation
\citep{constrained_decoding} realizes $\Phi$ via the pre-decoder
locus: the proof object dictates the admissible label set per step,
and a constrained decoder masks logits.

\begin{prop}[Subsumption, informal]
Procedural compilation $\subseteq \mathsf{Comp}(M)$ for $M \in
\mathsf{TLTS}_{\mathrm{func}}$. Ontological compilation $\subseteq
\mathsf{Comp}(M)$ for $M \in \mathsf{TLTS}_{\mathrm{rel}}$.
Procedural compilation is a strict specialization of ontological
compilation.
\end{prop}

The hull-indexing trick of procedural compilation works \emph{because}
$\delta$ is functional: when the answer is unique, geometric indexing
substitutes for soft attention. It does not lift to the relational
case in general but does lift to functional sub-fragments of an
Olog. This motivates a hybrid compilation strategy: in-FFN compile
the deterministic sub-Olog, pre-decoder mask the relational
remainder.

\section{Experiments}
\label{sec:expts}

\paragraph{What we are testing.} The framework makes architectural
predictions about constraint mechanics: that locus (B$'$) is
incomplete, that (C) saves latency on functional types, that the
cost of enforcement scales with prior--rule misalignment. These are
claims about the constraint geometry, not about whether a particular
trained LM happens to obey it. A synthetic prior holds model quality
fixed and varies only the constraint policy, so the resulting
numbers measure the mechanism. We return to a production attention
layer in Section~\ref{ssec:real} to check that the
synthetic-(B$'$) projection is a fair proxy for what a real
ontological mask does.

The harness simulates an unconstrained model as a hand-crafted
distribution $p_M(\ell \mid t)$; variants apply different
enforcement loci as sampling policies.\footnote{Code and full results in supplementary.
Variant (B) attention-layer masking is simulated by a sampling-time
projection $(B')$ admitting any label whose target is reachable from
the current state; the projection is faithful to the production
attention mask, validated in Section~\ref{ssec:real}.}

\paragraph{Setup.} We use a 7-type e-commerce Olog (Customer, Cart,
Item, Checkout, Payment, Order, Delivery) with one branching node
(Cart) and two terminals. We test two priors: \textsc{Good} (80\%
mass on admissible labels per type) and \textsc{Bad} (70\% mass on a
random favorite label per type, often inadmissible). Variants:
(A) standard sample from $p_M$; (B$'$) sample from $p_M$ masked by
type-reachability of label-target; (C) deterministic step on
functional types, (D) elsewhere; (D) sample from $p_M$ masked by
direct-edge admissibility from the current state.

\subsection{Soundness, fluency, latency}

\begin{table}[h]
\centering
\small
\begin{tabular}{@{}llrrrrr@{}}
\toprule
Prior & Variant & Soundness & $\log P/\text{traj}$ & Mean len &
KL/step & $\mu$s/step \\
\midrule
\textsc{Good} & A   &  48.1\% & $-2.54$ & 2.57 & 0.22 & 11.0 \\
\textsc{Good} & B$'$ & 61.5\% & $-2.20$ & 2.74 & 0.22 & 10.2 \\
\textsc{Good} & C   & 100.0\% & $-1.48$ & 3.51 & 0.22 &  6.0 \\
\textsc{Good} & D   & 100.0\% & $-1.47$ & 3.48 & 0.22 &  8.5 \\
\midrule
\textsc{Bad}  & A   &   4.2\% & $-1.95$ & 1.73 & 1.12 &  7.9 \\
\textsc{Bad}  & B$'$ &  4.3\% & $-1.92$ & 1.75 & 1.13 & 10.2 \\
\textsc{Bad}  & C   & 100.0\% & $-7.42$ & 3.51 & 1.92 &  5.5 \\
\textsc{Bad}  & D   & 100.0\% & $-7.33$ & 3.48 & 1.91 &  8.8 \\
\bottomrule
\end{tabular}
\caption{Variant comparison on the e-commerce TLTS. $N{=}1000$
trajectories, max length $12$, seed $42$. Soundness is the fraction
of trajectories whose every step is in $\delta$.}
\label{tab:variants}
\end{table}

Three findings (Table~\ref{tab:variants}). First, (C) and (D) reach
100\% soundness regardless of prior alignment; (A)'s soundness
tracks prior quality. Second, the fluency cost of enforcement scales
with prior--$\delta$ misalignment: under \textsc{Good}, $\log P$
barely shifts; under \textsc{Bad}, trajectories are $\approx 5.5$
nats less likely under the model's own distribution. KL-per-step is
a cheap runtime diagnostic of this cost. Third, (C) is $\approx
30\%$ faster per step because deterministic steps on functional
types (4 of 5 non-terminals here) skip the sampling pipeline.

\subsection{Reachability is not admissibility}
\label{ssec:reach}

(B$'$) is barely better than (A) on soundness in
Table~\ref{tab:variants}: $61.5\%$ vs.\ $48.1\%$ under \textsc{Good},
$4.3\%$ vs.\ $4.2\%$ under \textsc{Bad}. Reachability admits the
\emph{destination} (type $t'$ targeted by some label) without
requiring the \emph{step-edge} to exist as a $\delta$-triple from
the current state. The trajectory then breaks at the next step.

This is sharper at the limit. On a cyclic variant of the TLTS
(adding Delivery $\to$ Customer), reachability becomes 100\%
universal across type pairs and (B$'$) collapses:

\begin{table}[h]
\centering
\small
\begin{tabular}{@{}lrrr@{}}
\toprule
Topology & Reachability & (B$'$) snd, \textsc{Good} & (B$'$) snd, \textsc{Bad} \\
\midrule
Acyclic &   34.7\% &  65.0\% & 3.9\% \\
Cyclic  & 100.0\% &   8.0\% & 0.0\% \\
\bottomrule
\end{tabular}
\caption{(B$'$) collapses under universal reachability; (C)/(D)
remain 100\% on both topologies (omitted).}
\label{tab:cyclic}
\end{table}

\textbf{Implication.} Reachability masking at the attention layer
does not suffice; the decoder must also enforce direct-edge
admissibility. Production systems that enforce $\delta$ only at the
attention layer are unsound.

\subsection{Topology sweep: where the (C)/(D) tradeoff crosses}

We sweep a parameterized TLTS family (length-5 chain plus $k \in
\{0, \ldots, 5\}$ optional side-branches). The functional ratio is
$(5-k)/5$. The (C)/(D) latency ratio crosses 1.0 at functional ratio
$\approx 0.2$:

\begin{table}[h]
\centering
\small
\begin{tabular}{@{}rrrr@{}}
\toprule
$k$ & fn-ratio & (C)/(D) latency & (B$'$) snd, \textsc{Good} \\
\midrule
0 & 1.00 & \textbf{0.39} & 54.0\% \\
1 & 0.80 & 0.67 & 62.3\% \\
2 & 0.60 & 0.80 & 62.2\% \\
3 & 0.40 & 0.91 & 61.8\% \\
4 & 0.20 & 0.99 & 60.1\% \\
5 & 0.00 & 1.19 & 73.8\% \\
\bottomrule
\end{tabular}
\caption{Topology sweep. (C)/(D) crossover at functional ratio
$\approx 0.2$; (B$'$) never approaches 100\% regardless of topology.}
\label{tab:sweep}
\end{table}

Below the crossover, (C)'s deterministic-step shortcut on functional
types pays for its bookkeeping overhead; above it, the overhead
dominates. \textbf{Deployment heuristic}: use (C) when the
functional fragment exceeds $\approx 20\%$ of the non-terminal type
set; otherwise use (D).

\subsection{Validation against production code}
\label{ssec:real}

We bridge the cyclic TLTS into an OlogGraph and run the production
ontological attention layer's mask construction on a 6-token
trajectory window. The mask admits 6 self-pairs, 6 $\delta$-direct-edge
pairs, and \textbf{24 reachable-only pairs}---a 4:1 leakage ratio.
With random-init weights, forward-pass attention mass distributes
roughly uniformly across admitted pairs, landing 66.7\% on
reachable-only and only 16.4\% on $\delta$-direct-edge pairs. The
(B$'$) projection matches what the production layer does. The
production mask does not push mass toward direct-edge pairs;
soundness must come from the decoder.

\section{Verification Protocol and Audit Certificates}
\label{sec:verif}

A TLTS-compiled transformer admits a four-step post-hoc verifier
\citep{constrained_decoding}. Given the input, the forward-pass
trace $h_0, \ldots, h_n$, and the declared TLTS:

\begin{enumerate}
\item \textbf{Decode.} Apply $E^{-1}$ (or nearest-neighbor) to each
  $h_i$ to recover candidate type $\hat{t}_i$; reject if any
  residual falls outside the codomain of $E$ by more than
  $\epsilon$.
\item \textbf{Recover labels} from the token-emission log.
\item \textbf{Check transitions.} For each $i$, verify
  $(\hat{t}_{i-1}, \ell_i, \hat{t}_i) \in \delta$.
\item \textbf{Check mask coherence.} Inspect attention-weight
  matrices: confirm no mass exceeds threshold on $(q,k)$ pairs with
  $A(\hat{t}_q, \hat{t}_k) = 0$.
\end{enumerate}

\paragraph{Upstream gating changes the verifier's role.} Alone, the
verifier is flow-control: failed generations retry and some inputs
hang. Downstream of in-FFN or pre-decoder enforcement, it should
never fail except on numerical drift, becoming a monitor and a
public audit certificate.

\paragraph{Audit certificate.} We ship a reference implementation
that emits a JSON certificate per trajectory containing: a
fingerprint (SHA-256-prefix over sorted $T, L, \delta$) for drift
detection; per-step records ($t_{\text{in}}, \ell, t_{\text{out}}$)
plus diagnostics (in-$\delta$ flag, prior-argmax, forced flag,
masked-KL); and a summary. A separate verifier consumes the JSON
and the TLTS, returning pass/fail with reasons. It needs no model
weights, no PyTorch, and no GPU. The reference verifier passes
adversarial tests: tampering with a state-out (e.g., to a type
$\notin T$) is caught by the per-step $\delta$ lookup; TLTS drift is
caught by both fingerprint mismatch and per-step failures.

We recommend that compiled-transformer papers ship the JSON
certificate alongside outputs and the TLTS spec alongside weights,
with a verifier that runs without the model. Without this,
compiled-transformer claims cannot be independently checked.

\section{Related Work}
\label{sec:related}

\paragraph{Neurosymbolic AI.} This work sits in the third wave of
neurosymbolic integration \citep{garcez2023nesy3rd, deraedt2020nesy},
which integrates structure rather than couples modules. DeepProbLog \citep{manhaeve2018deepproblog} and Logic
Tensor Networks \citep{badreddine2022ltn} align learned
representations with logical structure via the loss function; we
work at the architectural level (attention masks, FFN gates).

\paragraph{Compiled transformers.} Tracr \citep{lindner2023tracr}
compiles RASP programs into transformers; ALTA generalizes the
target. Theoretical work establishes Turing-completeness
\citep{perez2021turing, bhattamishra2020computational}. We treat
these constructions as the functional-$\delta$ case of a broader
pattern.

\paragraph{Type-safe attention and constrained decoding.} The olog
formalism \citep{spivakkent2012ologs,spivak2014category} provides
the categorical substrate; the type-safe attention layer built on
it is reference code in our supplementary
\citep{constrained_decoding}. Constrained decoding
\citep{willard2023outlines, geng2023grammar} is the grammar case
of pre-decoder enforcement; we extend it to typed transition systems.

\paragraph{Verification.} Mechanistic interpretability
\citep{conmy2023acdc} inspects circuits in trained models; we
inspect forward passes against a declared compilation specification,
which is a different problem with different failure modes.

\section{Discussion}
\label{sec:disc}

\paragraph{Over-constraint and the precision--fluency tradeoff.}
Constraint enforcement costs fluency: tight $\delta$ buys soundness
at the price of likelihood. Known failure modes
(mass-redistribution drift, premature lock-in, confidence collapse,
dead-end states, topic-drift suppression) appear when $p_M$ and
$\delta$ disagree. Mitigations, in increasing order of effort:
enrich the Olog (most over-constraint is an under-specified
ontology), beam search with constraints, multi-proof rerank,
hierarchical staging, backoff and abstention, and constraint-aware
fine-tuning. Soft constraints are tempting but forfeit the
soundness guarantee.

\paragraph{Model-size implications.} Standard transformers devote
parameters to suppressing implausible domain compositions: the
distributional regularities of what does not happen in context.
TLTS-compilation moves that knowledge into the Olog and frees model
capacity for fluency and within-admissible-set selection. We predict
a shift in domain-specific deployment from weight-heavy with a thin
ontology to weight-light with a rich ontology, while open-domain
models still need capacity. We do not test this here;
constraint-aware fine-tuning is future work.

\paragraph{Limitations.} The functor account is informal; promoting
Proposition~\ref{prop:sound} to a theorem requires fixing
residual-stream geometry and gate-sharpness precision. The
empirical study uses synthetic priors; trained-model experiments
are future work. Verification scales as $|T|^2 |L|$, requiring lazy
evaluation for very large ontologies.

\section{Conclusion}

TLTS-compilation unifies two previously separate threads, type-safe
attention and procedural transformer compilation. It yields hybrid
architectures, a post-hoc verification protocol, and a JSON-based
audit certificate. All three are buildable today. We test the
main empirical prediction next: that domain-specific deployment can
substitute ontology engineering for raw scale.

\bibliography{references}

\end{document}
